\title{Controlling Text-to-Image Diffusion by Orthogonal Finetuning}

\begin{abstract}
State-of-the-art text-to-image diffusion models excel at synthesizing realistic images from descriptive textual inputs. Nevertheless, effectively steering these models toward a range of downstream tasks remains a critical unresolved issue. In response, we present Orthogonal Finetuning~(OFT), a systematic finetuning approach for customizing text-to-image diffusion models for various applications. OFT is distinct from previous approaches in that it rigorously maintains the hyperspherical energy, a measure of pairwise relationships between neurons on the unit hypersphere. Our analysis reveals that this conservation is essential for sustaining the underlying semantic generation proficiency of diffusion models. To enhance the robustness of finetuning, we also introduce Constrained Orthogonal Finetuning (COFT), which adds a supplementary radius constraint on the hypersphere. Concretely, we focus on two representative text-to-image finetuning scenarios: subject-driven generation—where the challenge lies in synthesizing images of a specific subject in novel contexts using a small number of reference images and a prompt—and controllable generation, where the aim is to incorporate external control cues. Our empirical evaluations demonstrate that OFT consistently surpasses prior methods in both generation fidelity and convergence efficiency.
\end{abstract}

\section{Introduction}

Modern text-to-image diffusion models~\cite{saharia2022photorealistic,ramesh2022hierarchical,rombach2022high} have set new standards in producing high-resolution images with nuanced, text-directed content. Yet, reliance solely on textual prompts can lead to imprecision and fails to offer the level of nuanced control required for more sophisticated image manipulation tasks. This motivates our focus on two key categories within text-to-image generation:

\begin{itemize}[leftmargin=*,nosep]

    \item \textbf{Subject-driven generation}~\cite{ruiz2023dreambooth}: In this setting, the objective is to create new renditions of a specific subject—identified from a handful of sample images—situated in varied scenarios, with generation guided by an accompanying text prompt.
    \item \textbf{Controllable generation}~\cite{zhang2023adding,mou2023t2i}: This task extends generative control through supplementary signals (\emph{e.g.}, edge maps, segmentation masks), ensuring that the synthesized images reflect both the control input and the text prompt.
\end{itemize}

At their core, both tasks revolve around the challenge of finetuning text-to-image diffusion models in a manner that retains the generative capabilities achieved during pretraining. We outline the essential qualities of an ideal finetuning approach as follows: (1) \emph{training efficiency}, which refers to minimizing the number of trainable parameters and reducing the total number of training epochs required, and (2) \emph{generalizability preservation}, ensuring that the model continues to produce high-fidelity and varied outputs. Typically, finetuning is performed by either making incremental updates to neuron weights using a small learning rate (\emph{e.g.}, \cite{ruiz2023dreambooth}), or by incorporating a lightweight module with its own set of re-parameterized weights (\emph{e.g.}, \cite{hulora2022,zhang2023adding}). However, these methods, while straightforward, do not offer guarantees that the model’s generative strengths developed during pretraining will be maintained. Furthermore, there is a lack of systematic insight into the development of robust finetuning strategies or the determination of key hyperparameters like training epochs. One central challenge is the absence of a clear metric for assessing how much of the pretrained model’s generative capabilities are preserved. Most current finetuning approaches work under the implicit assumption that a smaller Euclidean distance between the updated and original model weights signals better preservation of pretraining, and therefore, employ extremely conservative learning rates. Although this assumption might sometimes hold, we contend that relying solely on Euclidean distance does not adequately reflect the preservation of underlying semantics. Thus, adopting a more structurally informed metric to capture the divergence between the finetuned and original models could substantially enhance both the retention of pretraining performance and the overall stability of finetuning.

Drawing upon empirical findings that hyperspherical similarity serves as an effective proxy for capturing semantic relationships~\cite{liu2017deep,liu2018decoupled,chen2020angular}, our approach utilizes hyperspherical energy~\cite{liu2018learning} to encapsulate the relational dynamics among neurons. Here, hyperspherical energy is calculated as the aggregate of hyperspherical similarities (\emph{e.g.}, cosine similarities) for all neuron pairs within a single layer, thereby quantifying the uniformity of neuron distribution across the unit hypersphere~\cite{liu2021learning}. We posit that an optimally finetuned model should exhibit negligible variation in hyperspherical energy relative to its pretrained counterpart. A straightforward approach might involve incorporating a regularization term to maintain constant hyperspherical energy throughout the finetuning phase; however, this does not necessarily minimize the energy disparity effectively. Instead, we exploit the invariance characteristic of hyperspherical energy—namely, that pairwise hyperspherical similarity remains constant when the same orthogonal transformation is applied across all neurons. Building on this property, we introduce Orthogonal Finetuning~(OFT), which enables the adaptation of large-scale text-to-image diffusion models for downstream tasks while preserving their original hyperspherical energy. The foundational principle of OFT is to learn a layer-wise, shared orthogonal transformation of neuron representations, ensuring conservation of their pairwise angular relationships. This process may also be interpreted as redefining the canonical coordinate basis for the neurons in each layer. Recognizing that minimizing the Euclidean distance between finetuned and pretrained model weights aids in retaining pretrained capabilities, we additionally present a variant named Constrained Orthogonal Finetuning~(COFT). COFT enforces the constraint that the adjusted model resides within a hypersphere of fixed radius centered at the pretrained neuron representations.

The rationale for employing orthogonal transformations in the context of finetuning neurons draws in part from the principles of the 2D Fourier transform. In this framework, images are decomposed into their magnitude and phase spectra. Of particular importance is the phase spectrum, which encodes the angular relationship between the input and the basis and is chiefly responsible for retaining semantic content. Illustratively, if the phase spectrum of an image is combined with an arbitrary magnitude spectrum, the reconstructed image continues to reflect the original semantics. This observation underlines that altering the direction of neurons is central to effecting semantic adjustments in the output—an objective central to both subject-driven and controllable image generation. However, unconstrained modifications of neuron orientations can significantly undermine the generative performance acquired during pretraining. To manage this, we advocate for maintaining the angles between every pair of neurons, stemming from the premise that such angular relationships are critical to the knowledge encapsulated within neural networks. Accordingly, we introduce the idea of learning layer-specific orthogonal transformations for neurons, thereby ensuring the preservation of hyperspherical energy.

Further motivation is derived from the work on orthogonal over-parameterized training~\cite{liu2021orthogonal}, where classification neural networks are trained from initialization by applying orthogonal transformations. The rationale is that random initializations naturally result in networks with low expected hyperspherical energy, and \cite{liu2021orthogonal} aims to maintain this low energy throughout the training process—an objective supported by prior findings linking lower energy to improved classification generalization~\cite{liu2018learning,lin2020regularizing}. Notably, \cite{liu2021orthogonal} demonstrates that orthogonal transformations possess enough expressiveness to enable training of robust classifiers. Our approach, by contrast, is tailored towards finetuning text-to-image diffusion models with a dual focus on enhancing controllability and optimizing downstream generative outcomes. There are two key distinctions between OFT and the strategy employed by \cite{liu2021orthogonal}. Firstly, unlike \cite{liu2021orthogonal} which aims to decrease hyperspherical energy, our OFT approach seeks to maintain the original energy of the pretrained model, thereby safeguarding its structural integrity during finetuning—a crucial consideration in diffusion models, where excessive reduction of hyperspherical energy can compromise semantic representations. Secondly, OFT introduces constraints to minimize divergence from the pretrained model, resulting in a constrained version of the transformation, whereas \cite{liu2021orthogonal} operates without such restrictions. Finetuning, therefore, requires a careful balance between adaptability and retention of preexisting features, and we contend that OFT provides an effective means of achieving this equilibrium. The principal contributions of our work are summarized as follows:

\begin{itemize}[leftmargin=*,nosep]

    \item We introduce Orthogonal Finetuning, a new approach for enhancing the controllability of text-to-image diffusion models. Additionally, to reinforce finetuning stability, we present a constrained version that restricts the angular shift from the original pretrained parameters.
    \item In contrast with conventional finetuning strategies, OFT maintains the hyperspherical energy during model adaptation, a property we show—through empirical analysis—to be vital for retaining the semantic generation capabilities of the original model.
    \item Our method, OFT, is evaluated on two application domains: subject-driven generation and controllable generation. Extensive experiments confirm that OFT outperforms previous methodologies in terms of output quality, speed of convergence, and robustness of the finetuning process. Furthermore, OFT demonstrates superior data efficiency by achieving robust convergence with fewer training images and epochs.
    \item For assessing controllable generation, we develop a control consistency metric to quantitatively measure the alignment of generated images with control signals. This metric compares the estimated control cues in the synthesized output against the initial inputs, thereby providing evidence of OFT's strong capability for controllable generation.
\end{itemize}

\section{Related Work}

\textbf{Text-to-image diffusion models}. Significant advances~\cite{saharia2022photorealistic,ramesh2022hierarchical,rombach2022high,gu2022vector,nichol2022glide} have recently occurred in the domain of text-to-image synthesis, attributable to progress in diffusion-based generative models~\cite{ho2020denoising,song2021score,song2021denoising,dhariwal2021diffusion} and multimodal representation learning~\cite{su2020vlbert,lu2019vilbert,radford2021learning,wang2021simvlm,li2022blip,li2023blip,alayrac2022flamingo,schuhmann2022laion}. GLIDE~\cite{nichol2022glide} and Imagen~\cite{saharia2022photorealistic} implement diffusion in the pixel domain. GLIDE incorporates a jointly trained text encoder and diffusion prior using paired text-image datasets, while Imagen operates with a fixed pretrained text encoder. Alternatively, Stable Diffusion~\cite{rombach2022high} and DALL-E2~\cite{ramesh2022hierarchical} utilize latent diffusion modeling: Stable Diffusion employs VQ-GAN~\cite{esser2021taming} to construct its visual codebook for latent encoding, and DALL-E2 utilizes CLIP~\cite{radford2021learning} to create a shared latent representation unifying images and text. Beyond diffusion models, research has explored generative adversarial networks~\cite{reed2016generative,zhang2017stackgan,xu2018attngan,li2019controllable} and autoregressive models~\cite{ramesh2021zero,ding2021cogview,wu2022nuwa,yuscaling} for text-conditional image generation. OFT functions as a general finetuning paradigm compatible with any text-to-image diffusion architecture.

\textbf{Subject-driven generation}. Approaches such as \cite{avrahami2022blended,nichol2022glide} add a user-supplied mask as an auxiliary condition to avoid subject alterations. Inversion-based techniques~\cite{choi2021ilvr,dhariwal2021diffusion,ramesh2022hierarchical,gal2022image} enable modifying backgrounds or context without affecting the central subject. The framework in \cite{hertz2022prompt} allows both local and global content edits without the requirement for an explicit input mask. Despite these strategies, identity preservation often remains inadequate. In Pivotal Tuning~\cite{roich2022pivotal}, model parameters are adapted around an inverted latent vector, incorporating a regularization objective to maintain identity consistency. Similarly, \cite{nitzan2022mystyle} constructs a personalized face prior using multiple images of the same individual. \cite{casanova2021instance} supports the creation of different instance variants, yet sometimes at the expense of instance-specific traits. DreamBooth~\cite{ruiz2023dreambooth} introduces a bespoke token for each subject and, with a handful of reference images, applies finetuning using reconstruction and class-specific prior losses. While OFT is implemented following the DreamBooth setting, it replaces straightforward finetuning at a reduced learning rate with orthogonal updates to the model's weights.

\textbf{Controllable generation}. Image-to-image translation can be characterized as a problem of controllable image generation, where most earlier approaches rely on conditional generative adversarial networks~\cite{isola2017image,zhu2017toward,wang2018high,park2019semantic,choi2018stargan}. Recent research has shown that diffusion models are also effective in this setting~\cite{saharia2022palette,voynov2022sketch,wang2022pretraining}. More recently, the introduction of ControlNet~\cite{zhang2023adding} has enabled precise modulation of pretrained diffusion models via finetuning and incorporating additional control inputs, resulting in state-of-the-art controllable generation performance. Similarly, T2I-Adapter~\cite{mou2023t2i} finetunes pretrained diffusion models to improve control over image synthesis. Adopting the experimental framework of \cite{zhang2023adding,mou2023t2i}, we employ OFT to refine pretrained diffusion models, achieving superior controllability using less training data and a reduced number of tunable parameters. Importantly, OFT introduces no extra inference-time computation overhead.

\textbf{Model finetuning}. The practice of refining large pretrained models for specific downstream applications has become increasingly prevalent~\cite{devlin2018bert,brown2020language,he2022masked}. Within this domain, adaptation techniques (\emph{e.g.}, \cite{hulora2022,houlsby2019parameter,pfeiffer2020adapterfusion}) are extensively investigated in natural language processing. OFT is most closely related to LoRA~\cite{hulora2022}, which assumes an additive, low-rank adjustment to model weights during adaptation. In contrast, OFT modifies neuron weights via a multiplicative, layer-shared orthogonal transformation, which mathematically guarantees the preservation of pairwise angular relationships among neurons within each layer, thereby offering greater training stability.

\section{Orthogonal Finetuning}

\subsection{Why Does Orthogonal Transformation Make Sense?}

To motivate the use of orthogonal transformations for finetuning text-to-image diffusion models, we break down the issue into two points: (1) the rationale for adjusting neuron angles (\emph{i.e.}, their directions), and (2) the justification for choosing orthogonal transformation as the method for tuning these angles.

To address the first question, we build upon empirical findings from prior works \cite{liu2018decoupled,chen2020angular}, which suggest that differences in angular features are highly indicative of semantic distinctions. SphereNet~\cite{liu2017deep} demonstrates that constraining all network neurons to reside on a unit hypersphere during training maintains network capacity and can even enhance generalization, indicating that neuron orientations can encapsulate essential information from the data. To further highlight the role of neuron angles, we conduct a simple experiment illustrated in Figure~\ref{fig:toy_ae}, where a conventional convolutional autoencoder is trained on a dataset of flower images. During training, feature maps are generated via standard inner products ($z$ refers to the output element from the convolutional kernel $\bm{w}$ when applied to the input patch $\bm{x}$). In evaluation, we contrast three methods for obtaining the feature map: (a) employing the training-time inner product, (b) using the magnitude, and (c) extracting angular information. As presented in Figure~\ref{fig:toy_ae}, angular features alone are nearly sufficient to reconstruct the original images, whereas magnitude fails to provide informative content. It is important to note that the cosine-based (c) activation is not applied during training—training exclusively relies on inner products. These findings suggest that neuron orientations (angles) are the principal carriers of semantic content in the image inputs. Thus, to meaningfully alter the semantic representation in images, adjusting neuron directions is likely to yield more significant effects.

For the second question, an intuitive approach for refining neuron orientations is to rotate or reflect all neurons within a layer collectively, naturally introducing an orthogonal transformation. Although one could conceive of more adaptable angular adjustments (e.g., rotating each neuron individually), our findings indicate that orthogonal transformations provide an effective balance between expressiveness and structured regularization. Furthermore, evidence from \cite{liu2021orthogonal} suggests orthogonal transformations possess sufficient expressivity for neural network training. To illustrate the benefit of enforcing orthogonality, we conduct a controlled generation experiment following the ControlNet~\cite{zhang2023adding} setup, with outcomes displayed in Figure~\ref{fig:ortho}. The results reveal that our standard OFT, under the orthogonality constraint, exhibits stable training and achieves precise control after about 20 epochs. By contrast, removing the orthogonality requirement causes OFT to fail at both realistic image generation and achieving controllable outputs. This experiment highlights the necessity of orthogonality in OFT for reliable performance.

\subsection{General Framework}

Traditional finetuning methods generally employ gradient descent with a relatively low learning rate to adjust the parameters of a model or a subset of its layers. The use of a small learning rate effectively biases the optimization towards minimal divergence from the pretrained parameters, thus standard finetuning is approximately equivalent to solving an optimization problem that reduces $\thickmuskip=2mu \medmuskip=2mu \| \bm{M} - \bm{M}^0\|$, with $\bm{M}$ denoting the finetuned weights and $\bm{M}^0$ representing those from the pretrained model. While this soft regularization is intuitively reasonable, it may offer excessive flexibility, particularly when dealing with large-scale models. To impose a more restrictive structure, LoRA introduces a low-rank limitation on the updates such that $\thickmuskip=2mu \medmuskip=2mu \text{rank}(\bm{M} - \bm{M}^0)=r'$, where $r'$ is typically chosen to be a small value. In contrast, OFT imposes a condition based on maintaining pairwise similarities among neurons: $\thickmuskip=2mu \medmuskip=2mu \| \text{HE}(\bm{M})-\text{HE}(\bm{M}^0) \|=0$, where $\text{HE}(\cdot)$ corresponds to the hyperspherical energy of the weight matrix. 

To clarify this further, consider a fully connected layer denoted by $\thickmuskip=2mu \medmuskip=2mu \bm{W}=\{\bm{w}_1,\cdots,\bm{w}_n\}\in\mathbb{R}^{d\times n}$, where the vector $\bm{w}_i\in\mathbb{R}^{d}$ refers to the $i$-th neuron and $\bm{W}^0$ denotes the pretrained weights for this layer. The layer output $\thickmuskip=2mu \medmuskip=2mu \bm{z}\in\mathbb{R}^n$ is obtained via $\thickmuskip=2mu \medmuskip=2mu \bm{z}=\bm{W}^\top\bm{x}$, where $\bm{x}\in\mathbb{R}^d$ is the incoming activation. The OFT approach is thus interpreted as driving the difference in hyperspherical energy between the updated and original layers towards zero:

\begin{equation}\label{eq:oft_principle}
\footnotesize
\min_{\bm{W}} \left\| \text{HE}(\bm{W})-\text{HE}(\bm{W}^0)\right\|~~\Leftrightarrow~~\min_{\bm{W}} \bigg{\|} \sum_{i\neq j} \|\hat{\bm{w}}_i-\hat{\bm{w}}_j\|^{-1}-\sum_{i\neq j} \|\hat{\bm{w}}^0_i-\hat{\bm{w}}^0_j\|^{-1}\bigg{\|}
\end{equation}

where $\thickmuskip=2mu \medmuskip=2mu \hat{\bm{w}}_i:=\bm{w}_i/\|\bm{w}_i\|$ represents the normalized form of the $i$-th neuron, and the hyperspherical energy for the layer $\bm{W}$ is defined by $\thickmuskip=2mu \medmuskip=2mu \text{HE}(\bm{W}):= \sum_{i\neq j} \|\hat{\bm{w}}_i-\hat{\bm{w}}_j\|^{-1}$. Notably, the optimal value for Eq.~\eqref{eq:oft_principle} is zero, achievable whenever the weights $\bm{W}$ differ from $\bm{W}^0$ solely through a rotation or reflection, i.e., when $\thickmuskip=2mu \medmuskip=2mu \bm{W}=\bm{R}\bm{W}^0$ with $\thickmuskip=2mu \medmuskip=2mu \bm{R}\in\mathbb{R}^{d\times d}$ being an orthogonal matrix (where $\det(\bm{R})$ equal to $1$ or $-1$ designates a rotation or a reflection, respectively). This principle forms the core of OFT: rather than updating every parameter independently, OF

\begin{equation}\label{eq:oft_general}
    \footnotesize
    \bm{z}=\bm{W}^\top\bm{x}=(\bm{R}\cdot\bm{W}^0)^\top\bm{x},~~~\text{s.t.}~\bm{R}^\top\bm{R}=\bm{R}\bm{R}^\top=\bm{I}
\end{equation}

Here, $\bm{W}$ refers to the weight matrix obtained via OFT-finetuning, and $\bm{I}$ denotes the identity matrix. The operational mechanism of OFT is depicted in Figure~\ref{fig:block}. Analogous to the role of zero initialization in LoRA, it is necessary for OFT to begin fine-tuning precisely from the pretrained parameter $\bm{W}^0$. To facilitate this, we set the initial value of the orthogonal matrix $\bm{R}$ as the identity, thus ensuring the fine-tuned model starts from the pretrained state. The orthogonality of $\bm{R}$ is preserved through differentiable orthogonalization techniques, as discussed in \cite{liu2021orthogonal,lezcano2019cheap}. Further, we elaborate on efficient methods for maintaining orthogonality.

\subsection{Efficient Orthogonal Parameterization}\label{sect:eff_ortho}

Conventional orthogonalization procedures, such as the Gram-Schmidt algorithm, although differentiable, are computationally intensive and often impractical for large-scale applications~\cite{liu2021orthogonal}. For greater computational tractability, we employ Cayley parameterization to construct the orthogonal matrix. Specifically, the orthogonal matrix is parameterized as $\thickmuskip=2mu \medmuskip=2mu \bm{R}=(\bm{I}+\bm{Q})(I-\bm{Q})^{-1}$, with $\bm{Q}$ being a skew-symmetric matrix so that $\thickmuskip=2mu \medmuskip=2mu \bm{Q}=-\bm{Q}^\top$. This approach significantly enhances efficiency, albeit with the constraint that Cayley parameterization is restricted to orthogonal matrices of determinant $1$, that is, the special orthogonal group. However, empirical evidence suggests that this restriction does not negatively impact performance. Even with Cayley parameterization, representing $\bm{R}$ becomes parameter-heavy for large $d$. To overcome this challenge, we adopt a block-diagonal structure for $\bm{R}$, consisting of $r$ blocks. The parameterization is thus given by:

\begin{equation}
    \footnotesize
    \bm{R}=\text{diag}(\bm{R}_1,\bm{R}_2,\cdots,\bm{R}_r)=\begin{bmatrix}
    \bm{R}_1\in \text{O}(\frac{d}{r}) &  &\\
    &\ddots & \\
    & & \bm{R}_r\in \text{O}(\frac{d}{r})
    \end{bmatrix}\in \text{O}(d)
\end{equation}

Here, $\text{O}(d)$ refers to the group of all orthogonal matrices of dimension $d$. The matrices $\thickmuskip=2mu \medmuskip=2mu \bm{R}\in\mathbb{R}^{d\times d}$ and $\thickmuskip=2mu \medmuskip=2mu \bm{R}_i\in\mathbb{R}^{d/r\times d/r},\forall i$ are both orthogonal. If we set $\thickmuskip=2mu \medmuskip=2mu r=1$, the block-diagonal orthogonal matrix reduces to an unconstrained orthogonal matrix. An orthogonal matrix with dimensions $\thickmuskip=2mu \medmuskip=2mu d\times d$ possesses $\thickmuskip=2mu \medmuskip=2mu d(d-1)/2$ free parameters, corresponding to a computational complexity of $\mathcal{O}(d^2)$. For the case of an $r$-block diagonal orthogonal matrix, the parameter count is $\thickmuskip=2mu \medmuskip=2mu d(d/r-1)/2$, yielding a complexity of $\thickmuskip=2mu \medmuskip=2mu \mathcal{O}(d^2/r)$. Optionally, parameter sharing can be used across all blocks, i.e., by setting $\thickmuskip=2mu \medmuskip=2mu \bm{R}_i = \bm{R}_j$ for any $i\neq j$, which further reduces the parameter complexity to $\mathcal{O}(d^2/r^2)$. Importantly, even with these parameter reductions, the resultant matrix $\bm{R}$ retains orthogonality, so the property of hyperspherical energy is fully preserved.

Next, we evaluate the parameter efficiency of OFT in comparison to LoRA. In LoRA, with a low-rank parameter $r'$, the number of parameters to be trained is $\thickmuskip=2mu \medmuskip=2mu r'(d+n)$. Assuming both $r$ and $r'$ are proportional to the hidden size $d$ (for instance, $\thickmuskip=2mu \medmuskip=2mu r=r'=\alpha d$ for some fixed constant $0<\alpha\leq 1$), then LoRA's parameter complexity scales as $\mathcal{O}(d^2+dn)$, whereas OFT’s scales as $\mathcal{O}(d)$. For example, consider a weight matrix of size $\thickmuskip=2mu \medmuskip=2mu 128\times 128$: when $r'=8$, LoRA uses $2,048$ trainable parameters, while OFT employs $960$ trainable parameters at $r=8$, assuming no block sharing.

\subsection{Constrained Orthogonal Finetuning}

OFT’s capacity can be further restricted by enforcing that the finetuned model does not stray far from the original pretrained model. To achieve this, COFT utilizes the following forward computation:

\begin{equation}\label{eq:coft_general}
    \footnotesize
    \bm{z} = \bm{W}^\top \bm{x} = (\bm{R}\cdot\bm{W}^0)^\top \bm{x},~~~\text{s.t.}~\bm{R}^\top \bm{R} = \bm{R}\bm{R}^\top = \bm{I},~\norm{\bm{R}-\bm{I}}\leq \epsilon
\end{equation}

where both orthogonality and a constraint on the distance to the identity matrix, bounded by $\epsilon$, are enforced. The Cayley parameterization from Section~\ref{sect:eff_ortho} is suitable for maintaining orthogonality. However, implementing the $\epsilon$-bounded deviation within the Cayley parameterized form is challenging. To facilitate understanding of the Cayley transform, we employ the Neumann series to approximate $\thickmuskip=2mu \medmuskip=2mu \bm{R} = (\bm{I}+\bm{Q})(I-\bm{Q})^{-1}$, leading to $\thickmuskip=2mu \medmuskip=2mu \bm{R} \approx \bm{I} + 2\bm{Q} + \mathcal{O}(\bm{Q}^2)$, assuming that the Neumann

Since the series converges in the operator norm, it follows that the constraint $\thickmuskip=2mu \medmuskip=2mu\|\bm{R}-\bm{I}\|\leq\epsilon$ may be reformulated and imposed within the Cayley transform. This translates into the condition $\thickmuskip=2mu \medmuskip=2mu\|\bm{Q}-\bm{0}\|\leq\epsilon'$, with $\bm{0}$ representing the zero matrix and $\epsilon'$ denoting a separate error hyperparameter distinct from $\epsilon$. The new restriction on $\bm{Q}$ is conveniently handled through projected gradient descent. To ensure that the orthogonal matrix $\bm{R}$ starts as the identity, we initialize $\bm{Q}$ as the zero matrix. COFT thus unites two explicit forms of regularization: minimal hyperspherical energy deviation and explicitly limited divergence from the original pretrained parameters. Although existing finetuning techniques typically introduce the latter regularization only implicitly, COFT formulates it in a more direct manner. While OFT offers robust performance, empirical results show that COFT enhances the stability of finetuning, likely due to the explicit constraint on deviation. Figure~\ref{fig:eps_coft} demonstrates the role of $\epsilon$ in determining COFT performance: increasing $\epsilon$ allows the COFT-finetuned model to behave more like OFT, while decreasing $\epsilon$ makes the model increasingly resemble the pretrained text-to-image diffusion baseline.

\subsection{Re-scaled Orthogonal Finetuning}

As an enhancement to the standard OFT, we introduce an additional learnable scaling factor for each neuron. The rationale is that modifying neuron magnitudes does not affect the hyperspherical energy since this magnitude is normalized. Formally, the network computes the forward pass as $\thickmuskip=2mu \medmuskip=2mu\bm{z}=(\bm{R}\bm{W}^0\bm{D})^\top\bm{x}$\footnote{\textbf{Errata}: The NeurIPS camera ready version incorrectly states the forward pass for re-scaled OFT as $\thickmuskip=2mu \medmuskip=2mu\bm{z}=(\bm{D}\bm{R}\bm{W}^0)^\top\bm{x}$. The released implementation and the experimental results in Appendix~\ref{app:rescaled} remain valid.} where $\thickmuskip=2mu \medmuskip=2mu \bm{D}=\text{diag}(s_1,\cdots,s_n)\in\mathbb{R}^{n\times n}$ is a diagonal matrix with learnable positive entries $s_1,\cdots,s_n$. Contrary to OFT, where only $\bm{R}$ is optimized (as given in Eq.~\eqref{eq:oft_general}), both $\bm{R}$ and $\bm{D}$ are trainable parameters in the re-scaled variant. This adjustment modestly increases the parameter count while yielding greater model adaptability. To isolate the effect of orthogonal transformations, our main experiments employ the original OFT; nevertheless, the re-scaled version generally delivers stronger results (refer to Appendix~\ref{app:rescaled}).

\section{Intriguing Insights and Discussions}

\textbf{OFT accommodates diverse neural architectures}. The OFT technique, in principle, can be deployed for any neural network architecture. Within Transformer models, LoRA is customarily applied to attention matrices~\cite{hulora2022}. To ensure comparability, our experiments restrict OFT finetuning solely to attention weights, mirroring the usage of LoRA. Apart from standard dense layers, OFT also adapts naturally to convolutional layers. The block-diagonal configuration of $\bm{R}$ gains particular significance in the convolutional context—a contrast to LoRA. If the number of blocks equals the input channels, each block modifies a unique neuron channel, analogous to training depth-wise convolution filters~\cite{chollet2017xception}. If instead all blocks in $\bm{R}$ are identical, neurons are transformed through a common orthogonal matrix across all channels. We present preliminary experiments exploring OFT-based convolutional layer finetuning in Appendix~\ref{app:convolution}.

\textbf{Connection to LoRA}. LoRA addresses parameter efficiency by introducing a low-rank matrix addition, thereby ensuring that the pretrained weight matrix's core information remains mostly unaltered. In contrast, OFT imposes an orthogonal (full-rank) transformation on the pretrained weights, safeguarding the retention of pretraining features by restricting transformations to those that do not disrupt information content. OFT’s forward operation may be represented as $\thickmuskip=2mu \medmuskip=2mu \bm{z}=(\bm{R}\bm{W}^0)^\top\bm{x}=(\bm{W}^0+(\bm{R}-\bm{I})\bm{W}^0)^\top\bm{x}$, where the component $\thickmuskip=2mu \medmuskip=2mu (\bm{R}-\bm{I})\bm{W}^0$ plays a similar functional role to LoRA’s low-rank adjustment. Since $\bm{W}^0$ is conventionally full-rank, OFT likewise achieves a low-rank update when $\thickmuskip=2mu \medmuskip=2mu \bm{R}-\bm{I}$ itself is of low rank. In analogy to LoRA’s use of a rank parameter $r'$, OFT employs a block-diagonal size parameter $r$ to limit the total number of trainable weights. Crucially, LoRA and OFT achieve parameter efficiency via fundamentally different mechanisms: LoRA leverages low-rank structure, while OFT takes advantage of sparsity, specifically through block-diagonal orthogonal constraints.

\textbf{Why OFT converges faster}. As shown in Figure~\ref{fig:toy_ae}, modifying neural directions results in the most significant semantic changes, which is the precise focus of OFT. Furthermore, from an optimization standpoint, OFT’s procedure can be understood as refining neurons on a smooth hypersphere, leading to improved optimization dynamics and convergence. These observations are further substantiated by empirical evidence in \cite{liu2021orthogonal}.

\textbf{Why not minimize hyperspherical energy}. Distinct from the approach in \cite{liu2021orthogonal}, our method does not target minimizing hyperspherical energy. For classification tasks, it is advantageous for neurons to be non-redundant; however, achieving minimal hyperspherical energy spreads neurons uniformly across the hypersphere, a property that may undermine the retention of useful pretrained representations during finetuning, making such minimization counterproductive in this context.

\textbf{Balancing flexibility and regularization in finetuning}. Our analysis reveals a fundamental compromise between flexibility and regularity in finetuning strategies. While conventional finetuning methods offer maximal adaptability, they are prone to instability and frequent model degradation. OFT, despite its simplicity, effectively mediates this compromise by maintaining the pairwise angles between neurons, striking a favorable equilibrium. Furthermore, the block-diagonal structure functions as an enhanced regularizer for the orthogonal transformation matrix.

\textbf{Inference efficiency maintained}. In contrast to ControlNet, our OFT method incurs no inference-time computational overhead for the adapted model. During inference, the orthogonal matrix $\bm{R}$ that has been learned can be directly applied to the original weight matrix $\bm{W}^0$ via multiplication, resulting in an updated weight matrix $\thickmuskip=2mu \medmuskip=2mu \bm{W}=\bm{R}\bm{W}^0$. Therefore, the inference throughput remains consistent with that of the original, unadapted model.

\section{Experiments and Results}

\textbf{Experimental protocol}. For our empirical studies, we employ Stable Diffusion v1.5~\cite{rombach2022high} as the foundational text-to-image generation architecture. In all cases, image samples from various techniques are selected randomly to ensure impartiality. DreamBooth~\cite{ruiz2023dreambooth} serves as our reference for subject-driven image creation, whereas ControlNet~\cite{zhang2023adding} and T2I-Adapter~\cite{mou2023t2i} guide our protocols for controllable image generation tasks. To ensure equitable comparison against LoRA, we confine the application of OFT and COFT to identical model layers used for LoRA finetuning. Additional implementation specifics and expanded results are reported in Appendix~\ref{app:settings}.

\subsection{Subject-driven Generation}

\textbf{Configuration}. Baselines for this task include DreamBooth~\cite{ruiz2023dreambooth} and LoRA~\cite{hulora2022}. The training objectives across all compared methods adhere to the loss formulation used in DreamBooth. For DreamBooth and LoRA, training procedures and hyperparameters follow the original publications' recommendations and optimal settings. Supplementary analyses and further empirical evidence can be found in Appendix~\ref{app:settings},~\ref{app:coft_oft},~\ref{app:more_qualitative},~\ref{app:failure}.

\textbf{Finetuning stability and convergence}. To begin, we assess the stability and convergence rate of finetuning for DreamBooth, LoRA, OFT, and COFT. The outcomes, illustrated in Figure~\ref{fig:teaser} and Figure~\ref{fig:db_convergence}, reveal that both COFT and OFT exhibit stable performance during diffusion model finetuning. After completing 400 iterations, DreamBooth and OFT-based approaches demonstrate strong control, whereas LoRA fails to adequately maintain subject identity. As training proceeds to 2000 iterations, DreamBooth starts to generate degraded or collapsed outputs, and LoRA is unable to correctly render the yellow shirt, often producing yellow fur instead. In contrast, OFT and COFT continue to display reliable and robust image control at this stage. These findings confirm that our OFT framework supports rapid as well as stable convergence in subject-driven generation tasks. Notably, this robustness to the number of finetuning iterations greatly reduces the complexity associated with determining appropriate iteration counts for various subjects. With OFT and COFT, it is feasible to set a higher iteration number without intricate hyperparameter tuning. Specifically, for COFT, choosing a suitable $\epsilon$ simplifies the selection process for both the learning rate and the number of iterations.

\textbf{Quantitative comparison}. Adopting the evaluation strategy of \cite{ruiz2023dreambooth}, we present a quantitative analysis of subject fidelity (using DINO~\cite{caron2021emerging} and CLIP-I~\cite{radford2021learning}), adherence to text prompts (CLIP-T~\cite{radford2021learning}), and diversity of generated samples (LPIPS~\cite{zhang2018unreasonable}). CLIP-I involves calculating the mean pairwise cosine similarity between CLIP embeddings of the generated and real samples. The DINO metric operates similarly to CLIP-I, but relies on ViT S/16 DINO features instead. CLIP-T measures the mean cosine similarity between the embedding of the text prompt and those of the generated images. For diversity, average LPIPS cosine similarities are computed among images generated from identical prompts and subjects. According to Table~\ref{table:dreambooth_final}, COFT and OFT significantly surpass DreamBooth and LoRA in both DINO and CLIP-I metrics, indicating stronger subject fidelity, while also attaining comparable or improved results for prompt alignment and diversity. Each experimental configuration is repeated 30 times with different random seeds on the same pretrained model to limit variability. Collectively, these results demonstrate that the OFT framework not only improves convergence and stability, but also consistently yields superior final outcomes.

\textbf{Qualitative comparison}. To provide a clearer insight into the advantages offered by OFT, we present a selection of randomly chosen outputs for subject-driven image generation in Figure~\ref{fig:dreambooth_final}. To ensure fairness, all methods are applied to the same finetuned model for image generation, and no text prompts are specifically adjusted for any individual example. For each approach, we select the model that achieves the highest validation CLIP scores. The results illustrated in Figure~\ref{fig:dreambooth_final} indicate that both OFT and COFT maintain robust semantic subject fidelity, in contrast to LoRA, which frequently struggles to preserve the subject’s identity (as seen in the bowl example, where LoRA completely loses the subject). Additionally, OFT and COFT exhibit much more precise responsiveness to text-based control, whereas DreamBooth, although competent in retaining subject identity, frequently does not align the image with the input prompt (see, for example, the first row in the bowl case). This qualitative analysis reveals that OFT successfully achieves enhanced control over image generation while simultaneously preserving subject characteristics. Furthermore, the performance of OFT is relatively insensitive to the number of training iterations, allowing it to perform well even with a larger iteration count—something that neither DreamBooth nor LoRA handles effectively. Additional qualitative samples are available in Appendix~\ref{app:more_qualitative}. Results from a human evaluation, provided in Appendix~\ref{app:human}, further support the superiority of OFT.

\subsection{Controllable Generation}

\textbf{Settings}. ControlNet~\cite{zhang2023adding}, T2I-Adapter~\cite{mou2023t2i}, and LoRA~\cite{hulora2022} serve as our baseline methods. We evaluate three demanding controllable generation scenarios in the main text: Canny edge to image~(C2I) on the COCO dataset~\cite{lin2014microsoft}, segmentation map to image~(S2I) on the ADE20K dataset~\cite{zhou2017scene}, and landmark to face~(L2F) using CelebA-HQ~\cite{karras2018progressive,xia2021tedigan}. All approaches are employed to finetune Stable Diffusion~(SD) v1.5 for 20 epochs on these datasets. Expanded experimental results can be found in Appendix~\ref{app:more_qualitative}, \ref{app:more_control_task}, and \ref{app:failure}.

\textbf{Convergence}. To assess the convergence properties of ControlNet, T2I-Adapter, LoRA, and COFT on the L2F task, we conduct both quantitative and qualitative analyses. Quantitatively, we utilize the mean $\ell_2$ distance between the target face landmarks (control) and those generated by the models as our primary evaluation metric. Figure~\ref{fig:face_convergence} illustrates the trajectory of landmark errors as each model is fine-tuned across varying numbers of epochs. Our findings indicate that COFT and OFT converge markedly more rapidly compared to their counterparts. For instance, LoRA requires 20 epochs to achieve the performance level attained by OFT at just the 8-th epoch. It is important to highlight that both OFT and COFT are implemented with a number of trainable parameters comparable to LoRA (and significantly fewer than ControlNet), yet they reach convergence much more swiftly than the alternative approaches. Additionally, OFT's efficient convergence is supported by results depicted in Figure~\ref{fig:teaser}; in particular, the example on the right demonstrates that OFT is notably more data-efficient than both ControlNet and LoRA, as it converges satisfactorily using only 5\% of the ADE20K dataset. For qualitative comparison, our focus is on OFT, COFT, and ControlNet, given that ControlNet achieves landmark errors most comparable to ours. The outputs in Figure~\ref{fig:face_convergence_qual} reveal that both OFT and COFT demonstrate stable convergence, progressively aligning the synthesized face pose with the control landmarks over time. By contrast, ControlNet exhibits a sudden onset of controllability only after the 8-th epoch, which corroborates the abrupt convergence pattern reported by \cite{zhang2023adding}. This stability in the convergence process makes the OFT framework particularly user-friendly, as its training dynamics are more gradual and foreseeable, simplifying the process of hyperparameter optimization for OFT.

\textbf{Quantitative comparison}. To assess the efficacy of controllable generation, we propose a control consistency metric. This metric operates by extracting the control signal from the generated output and comparing it against the input control signal. Specifically, for the C2I task, we measure performance using IoU and F1 score; for S2I, we report mean IoU alongside mean and overall accuracy; and for L2F, we evaluate using the mean $\ell_2$ distance between the control landmarks and those predicted. Further elaboration on these metrics can be found in Appendix~\ref{app:settings}. In all evaluations, we adopt the optimal hyperparameter configurations for each method considered. The outcomes, summarized in Table~\ref{table:control_gen}, indicate that both OFT and COFT significantly surpass alternative approaches in terms of control precision and strength. Notably, adapter-based techniques such as T2I-Adapter and ControlNet are characterized by slower convergence and inferior end performance. While LoRA outperforms ControlNet on the S2I task, its effectiveness is diminished on C2I and L2F tasks. Overall, LoRA exhibits a comparatively low performance ceiling, even when thoroughly tuned. In contrast, the OFT framework demonstrates continuous improvement with an increasing number of trainable parameters, suggesting untapped potential. Importantly, we highlight our quantitative controllable generation assessment as a distinct contribution, as it offers an accurate measure of the finetuned models' control capability, within the accuracy limits of the employed segmentation or detection model.

\textbf{Qualitative comparison}. Additionally, we conduct a qualitative evaluation of OFT and COFT in relation to leading approaches such as ControlNet, T2I-Adapter, and LoRA. As illustrated by the samples in Figure~\ref{fig:control_final}, both OFT and COFT consistently deliver images with superior realism and fidelity, while ensuring precise controllability. In the S2I task, LoRA is unable to generate images that conform to the provided segmentation map, whereas ControlNet, OFT, and COFT demonstrate effective adherence to the segmentation inputs. OFT and COFT, in particular, outperform ControlNet by synthesizing images with higher fidelity, richer details, and more coherent geometric arrangements—all achieved with a significantly reduced parameter count. When considering the C2I task, OFT and COFT successfully generate images that are semantically aligned with rough Canny edge maps, outperforming both T2I-Adapter and LoRA. For the L2F task, our methods offer the most precise facial pose control, maintaining accuracy even in the presence of complex pose variations. Across these three control scenarios, OFT and COFT consistently produce images of higher qualitative value compared to contemporary baseline methods, highlighting the robustness of the OFT framework for controlled generation. Further qualitative examples encompassing the three main control tasks are available in Appendix~\ref{app:control_exp}. Additionally, Appendix~\ref{app:more_control_task} provides evidence of OFT’s capacity to excel in broader control settings, such as dense pose to human body, sketch to image, and depth to image tasks.

\section{Concluding Remarks and Open Problems}

Inspired by the critical role of angular relationships among neurons in shaping visual semantics, this work introduces a straightforward yet robust fine-tuning strategy—orthogonal finetuning—for steering text-to-image diffusion models. Our approach is particularly tailored to two primary tasks: subject-driven generation and controllable generation. In contrast to prior techniques, OFT offers improved controllability and more stable fine-tuning while reducing the number of trainable parameters required. Furthermore, OFT incurs no extra inference computational costs, facilitating practical and efficient deployment.

Several noteworthy open questions arise from our method. To begin with, OFT maintains orthogonality using the Cayley parametrization, which necessitates a matrix inversion—a step that somewhat constrains its scalability. Although a block diagonal parametrization is employed to partially address this bottleneck, finding a differentiable and faster matrix inversion technique remains an unresolved challenge. Additionally, OFT holds promise for compositionality, since orthogonal matrices learned from distinct OFT finetuning tasks can be composed via multiplication, with their product still being orthogonal. It remains to be explored whether these combined orthogonal matrices retain information from each of their respective downstream tasks. Lastly, the model’s parameter efficiency heavily relies on the imposed block diagonal structure, which can introduce specific biases and reduce flexibility. Developing approaches that enhance parameter efficiency with fewer constraints and biases is an important avenue for future research.

\end{document}